\section{D10 Heat-Kernel Calibration Supplement}
\label{app:particle-d10-heat-kernel}

This appendix carries the D10 compact-group heat-kernel items needed by the particle surface.
The theorem-bearing leaf for the same-overlap thermal/Casimir law sits on the
screen-microphysics paper surface at fixed cutoff. What is recorded here is the
compact-group / Peter--Weyl lift and beta-shift bookkeeping used by the particle-side
calibration branch. The public D10 status is explicitly Phase II rather than recovered-core
Phase I. In particular, this appendix does not prove that OPH derives the physical
one-loop beta coefficients from edge sectors alone.

\subsection{One-Loop Calibration Frame}

At one loop, if couplings unify at \((M_U,\alpha_U)\),
\[
\alpha_i^{-1}(M_Z)=\alpha_U^{-1}+\frac{b_i}{2\pi}\ln\frac{M_U}{M_Z}.
\]
On the D10 branch these are the branch values \((M_U(P),\alpha_U(P))\) emitted by the forward
D10 solve; they are not inferred by inverse readback from measured low-energy couplings.
Writing \(A_i:=\alpha_i^{-1}(M_Z)\) and \(L:=\ln(M_U/M_Z)\), one finds
\[
L=\frac{2\pi}{b_1-b_2}(A_1-A_2),
\]
and the corresponding consistency relation for the D10 lane is
\[
A_3^{\mathrm{pred}}
=
\frac{b_3-b_2}{b_1-b_2}A_1
+
\frac{b_1-b_3}{b_1-b_2}A_2.
\]
This appendix keeps that algebra on the page as a calibration relation, not as a theorem that OPH
derives a supersymmetric UV spectrum.

\subsection{Benchmark coefficient comparison}

For the standard one-loop benchmark coefficients,
\[
\begin{array}{c|c|c}
\text{Model} & (b_1,b_2,b_3) & \alpha_s(M_Z)^{\mathrm{pred}} \\
\hline
\text{SM} & (41/10,-19/6,-7) & \approx 0.071 \\
\text{MSSM} & (33/5,1,-3) & \approx 0.116 \\
\text{Observed} & & 0.1179\pm 0.0010
\end{array}
\]
At one loop, the printed MSSM-style coefficients reproduce the observed closure far better than
the plain SM coefficients. On this D10 branch this is a benchmark matched by the D10
calibration branch, not a theorem that OPH derives a supersymmetric UV spectrum.

\subsection{Heat-Kernel Sector Law}

Read this subsection as the compact-group merge boundary above the fixed-cutoff microphysics
theorem. MaxEnt plus the bi-invariant compact-group package yield the familiar
\(d_R e^{-tC_2(R)}\) law once the Peter--Weyl lift is supplied; they are not a second,
independent proof of the finite-cutoff same-overlap Casimir branch.

\begin{theorem}[Heat-kernel edge-sector weights]
Under MaxEnt with bi-invariant constraints on a compact Lie group \(G\), the sector probabilities
take heat-kernel form:
\[
p_R(t)\propto d_R\,e^{-tC_2(R)},
\]
where \(d_R=\dim R\) and \(C_2(R)\) is the quadratic Casimir.
\end{theorem}

\begin{lemma}[Bi-invariant operators]
If \(G=\prod_i G_i\) is a compact semisimple Lie group written as a product of compact simple
factors, then any bi-invariant second-order differential operator on \(G\) has the form
\[
D=c_0\mathbf 1-\sum_i c_i\Delta_{G_i},
\]
where \(\Delta_{G_i}\) is the Laplace--Beltrami operator on the factor \(G_i\).
\end{lemma}

\begin{remark}
Bi-invariance is equivalent to \(D\in Z(U(\mathfrak g))\). Linear terms vanish because there are
no invariant vectors in the adjoint representation, and the quadratic part is proportional to the
Killing form on each simple factor, giving one independent coefficient per factor. The Casimir
element then acts as \(-\Delta_G\) in the regular representation.
\end{remark}

\subsection{Peter--Weyl Multiplicity and Beta Shifts}

By Peter--Weyl decomposition, the effective refinement-limit edge representation space is
\[
L^2(G)\cong \bigoplus_R V_R\otimes V_R^*.
\]
Entanglement traces over one factor give multiplicity \(d_R\) in \(p_R\), but loops see both
factors, so the effective multiplicity entering the calibration branch is
\[
N_{\mathrm{eff}}(R)=d_R\,p_R.
\]
The one-loop beta shift from edge sectors is then
\[
\Delta b_a=\sum_R p_R\,d_R\,T_a(R),
\]
where \(T_a(R)\) is the Dynkin index for gauge factor \(a\). This is the exact mathematical point
on the D10 lane where edge-sector heat-kernel weights feed the calibration branch. Matching to
MSSM-like one-loop running is a declared Phase II benchmark, not a promoted theorem-level
MSSM spectrum claim.

The missing bridge to theorem-level unification closure is explicit. One requires a derived
refinement-limit identification of which transportable sectors actually contribute to the running
carrier, a derived statistics/grading rule rather than an imposed fermionic-grading restriction,
controlled threshold and decoupling conventions on that same carrier, a derived representation
content/truncation rule for the sectors retained in the comparison, and a proof that the effective
loop multiplicities entering \(\Delta b_a\) are exactly the ones used in this calibration package
rather than nearby alternatives. On the declared runtime surface the
RG/matching/threshold/scheme packet is work in progress: scheme lock, threshold map, beta
provenance, and interval composition.

At the branch unification diffusion parameter \(t_U(P)\approx 1.64\) selected by the forward D10
solve, the benchmark shift is
\[
\Delta b\approx (2.49,4.38,3.97)
\qquad
\text{vs MSSM}
\qquad
(2.50,4.17,4.00).
\]
With the additional D10 calibration assumption of a fermionic-grading restriction to
half-integer \(\mathrm{SU}(2)\) sectors, this sharpens to
\[
\Delta b\approx (2.50,4.17,3.97),
\]
which matches the MSSM-style one-loop benchmark at the percent level under that declared
calibration package. Equivalently, the benchmark ratio
\[
\frac{\Delta b_3}{\Delta b_2}\approx 0.91
\]
sits close to the MSSM comparison value \(0.96\).
